\section{Background}\label{sec:related}
\paragraph{Segregation models as spin systems.}
\citet{schelling1971} showed that mild local preferences produce macroscopic segregation. The connection to spin systems has been developed by \citet{vinkovic2006}, who cast Schelling dynamics as a lattice gas with an energy, by \citet{stauffer2007}, who relate it to the Ising model, by \citet{dallasta2008}, who study phase behaviour of Schelling dynamics in one and two dimensions, and by \citet{gauvin2009}, who map the phase diagram of a model with vacancies. In these works the structure of the model is read off Monte Carlo output. We study a minimal member of the family: two types, no vacancies, ferromagnetic nearest-neighbour coupling and conserved composition. Our subject is the sampler and an estimator applied to its output, not a new sociological model.

\paragraph{Exchange dynamics and kernel perturbations.}
The Metropolis rule \citep{metropolis1953} and Kawasaki's conserved-order-parameter dynamics \citep{kawasaki1966} are standard \citep{newman1999}. A sampler is correct when its kernel is reversible with respect to the target and irreducible \citep{levin2009}. How far a perturbed kernel moves the stationary law is controlled by the mixing properties of the unperturbed chain and is typically bounded as $\|\tilde\pi-\pi\|\le K(P)\|\tilde P-P\|$ \citep{mitrophanov2005}. We use this only qualitatively and measure the shift exactly where the state space permits.

\paragraph{Coexistence, droplets and finite-size effects.}
The square-lattice Ising model has an exact critical temperature \citep{onsager1944} and spontaneous magnetisation \citep{yang1952}, and through the lattice-gas correspondence these fix the coexistence curve in the composition--temperature plane \citep{rowlinson1982}. In a finite system at fixed composition the passage from a homogeneous state to coexistence is rounded and goes through droplet formation, displaced from the thermodynamic binodal \citep{binder1987,biskup2002}. A maximum of an energy-fluctuation estimate is therefore an estimator whose relation to the binodal has to be established, not assumed.
